\section{The Method}

Not every result here carries the same weight, so each one is labeled by what it actually shows. A claim is \textbf{derived} when the construction forces it with no measured number fed in, like the fractional electric charges of one generation. It is \textbf{predicted} when a number computed from the structure is compared against experiment and could have come out wrong, like the weak mixing angle at unification. It is \textbf{reproduced} when a known piece of physics, set to run on the substrate, behaves as it already does in the laboratory, like color confinement or the hydrogen spectrum. Three more labels mark the edges. \textbf{Free} is a Standard Model number the model does not fix. \textbf{Open} is an unsolved problem with a named obstruction. \textbf{Posited} is the single experiential step set aside in the last section.

Each computational result also carries a depth grade from zero to three, and the test suite enforces it. Grade zero means the answer was put in by hand, so it shows nothing on its own. Grade one confirms a known mathematical fact, like the cell's directions forming the binary tetrahedral group. Grade two reproduces a known physical construction, like a lattice gauge theory. Grade three is the one that counts, a single base rule producing a measured result, checked against a control where the answer should come out no. Most results in a young program are grade one or two, which is stated rather than hidden, and any grade-three claim without a control is marked down, so no result can promote itself by assertion.

There is no randomness in the base, in the law or in the starting conditions. The Bell violation, the Born rule, and the genesis all run on fixed inputs, and a result is made robust by growing the system rather than by averaging over random seeds. A deterministic base that still yields the Born rule and the Tsirelson bound is a far stronger result than one that quietly feeds in quantum randomness by hand. The base is also discrete through and through, ternary values, a finite set of directions, a discrete time step, a rule that is a permutation, a finite symmetry group. Real numbers appear only as measured outputs or comparison data, never at the base, and the smooth continuum is what the fine discrete fabric looks like from far off. These are hard constraints, and they are what keep the model from cheating. Breaking any one is a cheat, not a refinement, and a phenomenon that does not emerge under them is recorded as a negative, never patched by adding a ninth element.

\begin{table}[ht]
\centering
\begin{tabular}{@{}lp{0.6\textwidth}@{}}
\toprule
constraint & what it forbids \\
\midrule
discrete & real numbers, infinite precision, a continuum in the base \\
deterministic & randomness, hidden seeds, anything but a fixed function \\
reversible & a built-in arrow of time in the closed bulk \\
charge-conserving & creating net charge from nothing \\
local & any action at a distance \\
no hidden state & memory or flags beyond the tones \\
one substance & a second kind of stuff \\
emergence gap & calling raw tones, or their first clusters, physics \\
\bottomrule
\end{tabular}
\caption{The constraints every base element must satisfy.}
\label{tab:constraints}
\end{table}

The model is not only argued, it is run. A finite, discrete simulator builds the exact mesh, runs the one rule over it beat by beat, and measures what emerges, so each question is a concrete experiment that either works or does not. The whole program is open research, and the full source code for the simulator and every experiment is available \cite{cluesurf2026repo}. That code was written with substantial AI assistance, which changes nothing about whether to trust it, since every result is deterministic and reproducible from the open suite, and reproducibility, not authorship, is the guarantee. Of the experiments in the suite, most are reproductions and confirmations, a smaller set are genuinely novel measured results, and a real fraction fail or stay open, concentrated in a few hard areas named as they arise, gravity, the sharp light cone, emergent dimension, and stable solitons most of all. The build fails only on a code crash or a broken consistency check, never on a scientific negative, because a negative is a result, not a bug.
